\section{Restarts}

This is the last thing we will implement in our SAT solver (for now, maybe). 
Restarts are exactly what it sounds like. We completely delete the trail and 
all notion of assignment. Why would we do that? A valid question. The idea is that 
perhaps while we've been solving, we got stuck in a bad part of the search space
and are just getting contradiction after contradiction. Rather than staying stuck there,
perhaps we could restart our solver so that it will end up in a different part of the 
search space that might be better. 

Makes enough sense, but how do we ensure that we don't just go back to the same 
part of the search space again? Well luckily because we've been learning, our solver 
isn't exactly the same as it initially was. It ``knows'' more, and since knowledge is 
power, the new learned clauses can force certain assignments and propagations earlier 
and that will nudge or push us into the right direction. 

So implement this, it's really really simple. All we need to do is after a certain number 
of conflicts, clear ONLY the trail. We don't clear out our VSIDS, clause measurements, 
reset the watched literals, nothing, nada, nahin. Why? Because the idea of restarts is 
that when the going gets tough, we get going. Notice how I said that we get going, not 
that we get forgetful? We go back to the beginning, remembering everything we've learnt, 
and use that to (hopefully) not repeat mistakes. If we forgot things like our VSIDS, we 
would make the same assignment with the exact same decision, and we'd have to rely only 
on our learned clauses to sway our decision making. But since VSIDS tells us which 
variables have been important in conflicts, it tells us which variables are best to start 
out with on our new try.

For our solver, I chose to restart every time we clear out the clause database. If you've 
been loosy-goosy about following my implementation or have been doing your own experiments
(nice), you can define your own condition. Just remember these things when you're doing a 
restart:
\begin{itemize}
    \item Clear the entire trail.
    \item The trail\_starts array must always have a decision level of 0.
    \item Clearing out var\_to\_trail is not strictly necessary as long 
            as your implementation is correct.
    \item If you do clear out var\_to\_trail, make sure it's resized properly.
    \item Reset qHead.
    \item Reset the decision level to 0.
    \item Unassign every variable.
    \item You might want to do a root level propagation.
\end{itemize}

Here are the results:
\begin{FVerbatim}
On all 1000 equations in uf20-91:

Restarts:                   385ms (~0.39 seconds).
Clause VSIDS:               347ms (~0.35 seconds).
M-size (10\%):              320ms (~0.32 seconds).

AMA (Y = 0.75 or 0.99):     274ms (~0.27 seconds).
EMA (X = 1, Y = 0.95):      266ms (~0.27 seconds).

Improved CDCL:              284ms (~0.28 seconds).
First UIPs:                 406ms (~0.41 seconds).
Backjumping:                380ms (~0.38 seconds).
Clause Learning:            425ms (~0.43 seconds).
Iterative DPLL (AbP):       328ms (~0.33 seconds).
Iterative DPLL (PbA):       324ms (~0.32 seconds).
Watched Literals:           349ms (~0.35 seconds).
Pure Literal Elimination:   1548ms (~1.55 seconds).
Unit Propagation:           868ms (~0.87 seconds).
Recursive Backtracking:     1,540,613ms (~25.68 minutes).
Brute Force:                1,770,851ms (~29.51 minutes).
Short-Circuited (Basic):    49,874ms (~50 seconds).
Parallel:                   10,589ms (~10.5 seconds).

On all 100 equations in Flat200-479:

Restarts:                   9,428ms  (~9.43 seconds).
Clause VSIDS:               19,181ms (~19.18 seconds).
M-size (10\%):              17,716ms (~17.72 seconds).
AMA:                        24,238ms (~24.24 seconds).

On all 1000 equations in UUF50.218.1000:

Restarts:                   1,824ms (~1.82 seconds).
Clause VSIDS:               1,804ms (~1.80 seconds).
M-size (10\%)               1,768ms (~1.77 seconds).

AMA:                        1,629ms (~1.63 seconds).
EMA (X = 1, Y = 0.95):      1,504ms (~1.50 seconds).

Improved CDCL:              2,385ms (~2.39 seconds).
First UIPs:                 14,012ms (~14.01 seconds).
Backjumping:                11,291ms (~11.29 seconds).
Clause Learning:            9,652ms (~9.65 seconds).
Iterative DPLL (AbP):       6,369ms (~6.37 seconds).
Iterative DPLL (PbA):       6,365ms (~6.37 seconds).
Watched Literals:           7,130ms (~7.13 seconds).
Pure Literal Elimination:   231,107ms (~3.85 minutes).
Unit Propagation:           140,041ms (~2.34 minutes).
\end{FVerbatim}
\noindent
Even though our solver is getting slower and slower on uf20-91, I think it's
actually really good that for the larger problems, it's getting substantially 
faster. Even though we want our algorithms and programs to work really quickly 
on small inputs, I think it's important for us to at some point prioritize the 
larger inputs. And given how most SAT problems are even larger than these, 
I think that's a good thing. 

\subsection{Assignment Polarity}

Back in the day, we had this thing called ``Pure Literal Elimination,'' but when 
we transitioned to CDCL, it went away \textit{forever}. Well dear reader, I'm here 
to tell you that it's still going to be gone forever, but its spirit lives on with us.
And by that, I mean, we're going to take the idea behind pure literal elimination and 
apply it to our restarts. 

You see, as we continue learning new clauses, we also encounter new literals. Earlier, 
back when we switched from brute force to recursive backtracking, we noticed that solving 
times decreased. Other than the main structure of the solvers, the only thing that changed
was in brute force, we started out with most variable assignments being False, while in 
recursive backtracking and beyond, they start out as True. I made a claim that this speedup 
was because there were less negative literals than there were positive literals. I won't 
say that I knew for sure, but if we stop for a second to think about it, it makes sense. 

When there are more positive literals, the likelihood of a variable's final assignment 
being True is higher than when there are more negative literals. This was the idea behind 
pure literal elimination, to an extent. Pure literal elimination worked because when a
literal only shows up with one polarity, if there is a solution, there is a 100\% chance 
that there is at least one solution where that pure literal's assignment matches its
polarity (T for positive, F for negative). We want to bring this idea into our assignment 
polarity. While we can't guarantee that literals will remain pure because of clause learning, 
we can at least make better decisions based on the current number of positive/negative occurrences 
of a variable.

To do this, we'll add in a \textit{polarity} array to our solver:
\begin{lstlisting}
    public interface:

    private interface:
        ...
        int[] polarity
        ...
\end{lstlisting}
\noindent
Instead of using a count for both positive and negative literals, instead, we'll 
use one counter for each variable. When the count is negative, we know that there 
are more negative occurrences of the variable, and when the count is nonnegative, we 
will pretend that are more positive occurrences of the variable (even though when the 
count is 0, that means it's balanced).

Because we're using one counter per variable, it might be tempting to use a literal's 
value directly in the count. After all, if we encounter one $+4$ and one $-4$, they'd 
cancel out to 0. Although this is logically correct, I'm worried about integer overflow 
for equations with millions of variables
\footnote{$2^32 \approx 4.3$ billion, so if we had one million variables, we could only 
handle ~43 occurrences of one polarity, which is pretty low, all things considered.}

Instead we'll settle for an ``if negative => count - 1, else count + 1'' kind of deal.
We will gather the counts in the setup phase of our solver, and we will update the counts 
only when we learn a new clause. However, it would be interesting to update the counts for 
all variables involved in a conflict, and see how that affects things. Because of how 
easy this code is to add as well, I'm gonna let you figure it out. Lastly, when we 
pick a variable via its VSIDS, we assign it based on its polarity. Also very easy to 
figure out, you got this, I believe in you.

Here are the results:
\begin{FVerbatim}
On all 1000 equations in uf20-91:

Polarity:                   369ms (~0.37 seconds).
Restarts:                   385ms (~0.39 seconds).
Clause VSIDS:               347ms (~0.35 seconds).
M-size (10\%):              320ms (~0.32 seconds).

AMA (Y = 0.75 or 0.99):     274ms (~0.27 seconds).
EMA (X = 1, Y = 0.95):      266ms (~0.27 seconds).

Improved CDCL:              284ms (~0.28 seconds).
First UIPs:                 406ms (~0.41 seconds).
Backjumping:                380ms (~0.38 seconds).
Clause Learning:            425ms (~0.43 seconds).
Iterative DPLL (AbP):       328ms (~0.33 seconds).
Iterative DPLL (PbA):       324ms (~0.32 seconds).
Watched Literals:           349ms (~0.35 seconds).
Pure Literal Elimination:   1548ms (~1.55 seconds).
Unit Propagation:           868ms (~0.87 seconds).
Recursive Backtracking:     1,540,613ms (~25.68 minutes).
Brute Force:                1,770,851ms (~29.51 minutes).
Short-Circuited (Basic):    49,874ms (~50 seconds).
Parallel:                   10,589ms (~10.5 seconds).

On all 100 equations in Flat200-479:

Polarity:                   8,580ms  (~8.58 seconds).
Restarts:                   9,428ms  (~9.43 seconds).
Clause VSIDS:               19,181ms (~19.18 seconds).
M-size (10\%):              17,716ms (~17.72 seconds).
AMA:                        24,238ms (~24.24 seconds).

On all 1000 equations in UUF50.218.1000:

Polarity:                   1,897ms (~1.90 seconds).
Restarts:                   1,824ms (~1.82 seconds).
Clause VSIDS:               1,804ms (~1.80 seconds).
M-size (10\%)               1,768ms (~1.77 seconds).

AMA:                        1,629ms (~1.63 seconds).
EMA (X = 1, Y = 0.95):      1,504ms (~1.50 seconds).

Improved CDCL:              2,385ms (~2.39 seconds).
First UIPs:                 14,012ms (~14.01 seconds).
Backjumping:                11,291ms (~11.29 seconds).
Clause Learning:            9,652ms (~9.65 seconds).
Iterative DPLL (AbP):       6,369ms (~6.37 seconds).
Iterative DPLL (PbA):       6,365ms (~6.37 seconds).
Watched Literals:           7,130ms (~7.13 seconds).
Pure Literal Elimination:   231,107ms (~3.85 minutes).
Unit Propagation:           140,041ms (~2.34 minutes).
\end{FVerbatim}
\noindent
That's a pretty decent speedup on the larger equations. Not sure how I feel 
about the speedup on UUF50.218.1000, but since it's not even 100ms across all 
1000 cases, I think I'll take it.